% Guide étudiant d'EP Manager (mod_ep).
% Compilation : latexmk -pdf guide-etudiant.tex
\documentclass[11pt,a4paper]{article}

% Version du guide : à mettre à jour à chaque modification, avec une ligne dans l'historique.
\newcommand{\titreguide}{Guide étudiant}
\newcommand{\versionguide}{1.1}
\newcommand{\dateguide}{novembre 2026}
\newcommand{\versionep}{1.0.2}
\newcommand{\versionsynthese}{1.0.2}

\usepackage{guide-commun}

\begin{document}

\pagedegarde{Suivre ses enseignements personnalisés, s'inscrire, déclarer, obtenir son
supplément au diplôme}{etu-tableau-de-bord}

\tableofcontents

\section*{Historique des versions}
\addcontentsline{toc}{section}{Historique des versions}

\begin{center}
\begin{tabularx}{\linewidth}{@{}l l X@{}}
\toprule
\textbf{Version} & \textbf{Date} & \textbf{Modifications} \\
\midrule
1.0 & novembre 2026 & Première version, pour EP Manager 1.0. \\
1.1 & novembre 2026 & Inscription directe depuis le catalogue (EP Manager 1.0.2). \\
\bottomrule
\end{tabularx}
\end{center}

\newpage

% ===========================================================================
\section{Votre espace d'enseignement personnalisé}

Vos enseignements personnalisés (EP) sont suivis dans le cours Moodle de votre promotion, du
début à la fin de votre cursus. Ouvrez l'activité \og Enseignement personnalisé\fg{} de ce
cours : vous arrivez sur votre tableau de bord.

\capture{etu-tableau-de-bord}{Tableau de bord de l'étudiant : ECTS retenus par année et par
type, face aux minimums exigés.}

Le tableau de bord montre :
\begin{itemize}
  \item pour chaque \textbf{année d'étude}, les ECTS retenus face au minimum exigé, et l'état
    \statut{Atteint} ou \statut{À compléter} ;
  \item le total sur le cursus, face au minimum de cursus s'il y en a un ;
  \item le détail \textbf{par type} d'EP (académique interne, stage, engagement étudiant,
    expérience professionnelle, sport, académique externe\dots), avec les plafonds ;
  \item la liste de \textbf{tous vos EP}, avec leur état et les décisions prises.
\end{itemize}

La barre de liens en haut de page mène au \chemin{Catalogue des EP}, à \chemin{Déclarer un
EP} et à \chemin{Mon supplément au diplôme}.

\begin{astuce}
Seuls les ECTS \textbf{validés} comptent. Au-delà du plafond d'un type, les ECTS validés restent
acquis mais ne sont plus comptés : ils apparaissent comme \og non retenus\fg{} et reviendraient
au décompte si le plafond était relevé.
\end{astuce}

% ===========================================================================
\section{S'inscrire à un EP du catalogue}

\subsection{Choisir un EP}

Le \chemin{Catalogue des EP} liste les EP proposés à votre promotion : intitulé, type, ECTS,
années concernées, places restantes, responsables, période d'inscription. Le lien
\bouton{Syllabus} ouvre la fiche de l'EP : objectifs, acquis d'apprentissage, contenu, méthodes,
conditions d'examen, programme des séances avec leurs intervenants, volume horaire. La fiche
se télécharge aussi en PDF.

\capture{etu-catalogue}{Catalogue des EP, vu par un étudiant.}

\capture{etu-syllabus}{Fiche du syllabus d'un EP.}

Le bouton \bouton{S'inscrire} n'apparaît que si l'inscription vous est possible. Sinon, la
colonne indique pourquoi :
\begin{itemize}
  \item \statut{Inscriptions fermées} : hors de la période d'inscription, ou EP clos ;
  \item \statut{Complet} : plus de place ;
  \item l'EP n'est pas ouvert à votre année d'étude ;
  \item \statut{Prérequis manquants} : l'EP exige d'avoir validé un autre EP, ou un nombre
    minimum d'ECTS. Si ses prérequis sont \textbf{bloquants}, l'inscription est refusée ; sinon,
    elle reste possible et le responsable en est averti.
\end{itemize}

\subsection{S'inscrire}

Sur la ligne de l'EP choisi, sélectionnez l'année d'étude au titre de laquelle il comptera (votre
année courante par défaut), puis cliquez sur \bouton{S'inscrire}. Votre inscription part
\textbf{en attente} : le responsable de l'EP la validera, avec une note si l'EP est noté, une fois
l'EP suivi. Vous restez sur la même page du catalogue, pour vous inscrire à d'autres EP.

Tant qu'elle est en attente, vous pouvez retirer votre inscription depuis votre tableau de bord
(\bouton{Retirer cet EP}) : la place est rendue.

% ===========================================================================
\section{Déclarer un EP hors catalogue}

Les activités suivies hors du catalogue -- engagement étudiant, expérience professionnelle,
sport, enseignement suivi dans un autre établissement\dots{} -- se déclarent depuis
\chemin{Déclarer un EP} :
\begin{enumerate}
  \item choisissez le \textbf{type} : sa consigne s'affiche, avec son plafond d'ECTS ;
  \item donnez un \textbf{intitulé}, l'\textbf{année d'étude}, le nombre d'\textbf{ECTS demandés}
    et une description de ce que vous avez fait ;
  \item joignez vos \textbf{justificatifs} (attestation, convention, diplôme\dots) : ils peuvent
    être obligatoires pour certains types ;
  \item enregistrez.
\end{enumerate}

\capture{etu-declarer}{Déclaration d'un EP hors catalogue.}

La page indique à qui la demande est transmise : selon le type, votre enseignant référent, la
DEVE, ou l'un puis l'autre. Votre enseignant référent est celui qui vous est attribué dans
l'activité \og Gestion des stages\fg{} du même cours.

Tant que personne ne l'a tranchée, vous pouvez modifier ou retirer votre demande. Une fois
tranchée, elle ne se modifie plus : en cas d'erreur, adressez-vous à la DEVE.

% ===========================================================================
\section{Suivre vos demandes}

Chaque EP de votre tableau de bord porte son état :
\begin{itemize}
  \item \statut{En attente}, avec la personne qui doit maintenant décider (responsable de l'EP,
    enseignant référent ou DEVE) : une demande peut passer par plusieurs étapes ;
  \item \statut{Validé}, avec les ECTS retenus et, pour un EP noté, la note ;
  \item \statut{Non validé}, avec le commentaire de la personne qui a décidé ;
  \item \statut{Retiré}.
\end{itemize}
Le détail d'un EP montre l'historique de toutes les décisions : qui, quand, à quel titre, avec
quel commentaire.

Vous recevez un message (courriel ou notification, selon vos préférences Moodle) à chaque
décision finale.

Les EP de type \textbf{stage} sont attribués automatiquement d'après vos stages complémentaires
validés par la DEVE dans \og Gestion des stages\fg{} : vous n'avez rien à déclarer.

% ===========================================================================
\section{Votre supplément au diplôme}

\subsection{Ce qu'il contient}

Le supplément au diplôme, au modèle européen Europass, décrit votre diplôme et votre parcours.
La page \chemin{Mon supplément au diplôme} montre :
\begin{itemize}
  \item les EP validés qui y figureront (rubrique~4.3), avec leurs ECTS et leur note ;
  \item vos \textbf{éléments complémentaires} (rubrique~6.1) : responsabilités associatives,
    prix, mobilités, certifications\dots{} que vous avez demandés et qui ont été validés ;
  \item le bouton \bouton{Aperçu (non certifié)}, qui produit le document en français ou en
    anglais, marqué \og APERÇU -- NON CERTIFIÉ\fg{} ;
  \item les suppléments \textbf{émis} par la DEVE, à télécharger.
\end{itemize}

\capture{etu-supplement}{Page \og Mon supplément au diplôme\fg{}.}

\subsection{Demander l'ajout d'un élément}

Le bouton \bouton{Demander l'ajout d'un élément} ouvre un formulaire :
\begin{enumerate}
  \item la \textbf{catégorie} (définie par la DEVE) ;
  \item la \textbf{période} (début, fin) ;
  \item l'\textbf{intitulé}, en français et, si possible, en anglais (pour le supplément édité
    en anglais) ;
  \item une description ;
  \item les \textbf{preuves} en PDF (attestation, certificat\dots), facultatives ou obligatoires
    selon le choix de la DEVE.
\end{enumerate}

\capture{etu-supplement-element}{Demande d'un élément complémentaire du supplément.}

La demande est examinée selon le circuit choisi par la DEVE (par exemple : votre enseignant
référent, puis la DEVE). Tant qu'elle est en attente, vous pouvez la modifier ou la retirer.
Vos preuves sont conservées par l'école avec la décision.

\subsection{Le supplément certifié}

La DEVE émet votre supplément en fin de cursus. Le document porte un \textbf{numéro} et un
\textbf{code de vérification} : toute personne à qui vous le présentez peut vérifier, sur la
page de vérification de l'école, qu'il a bien été émis et n'a pas été modifié. Si la DEVE émet
une nouvelle version, l'ancienne est signalée comme remplacée.

\end{document}
